\documentclass[conference]{IEEEtran}
\usepackage[utf8]{inputenc}
\usepackage{geometry}
\geometry{margin=1in}
\usepackage{graphicx}
\usepackage{hyperref}
\usepackage{caption}
\usepackage{subcaption}
\usepackage{listings}
\lstset{
  basicstyle=\tiny\ttfamily,
  language=Python
}
\usepackage{url}

\title{Local RAG-Based Chatbot for University Program Information}

\author{
    \IEEEauthorblockN{Student Name}
    \IEEEauthorblockA{Department of Computer Science\\University of Alberta\\Edmonton, Alberta, Canada}
}

\begin{document}

\maketitle

\begin{abstract}
This paper presents a local Retrieval-Augmented Generation (RAG) chatbot system for providing accessible information about university mathematics and statistics programs. Unlike cloud-based solutions that require API calls and raise privacy concerns, our system runs entirely locally using open-source tools including Ollama, LangChain, and Chroma. The chatbot scrapes university web pages, processes the content into structured JSON, creates a vector database for semantic search, and generates responses using the qwen3 language model. Our proof-of-concept demonstrates that a free, privacy-preserving chatbot can effectively answer student questions about program requirements, courses, and degree options. This work contributes to the growing body of research on local, offline-capable AI systems for educational purposes.
\end{abstract}

\IEEEkeywords
RAG, chatbot, local LLM, Ollama, education, university information system

% =============================================================================
% 1. INTRODUCTION
% =============================================================================
\section{Introduction}
\label{sec:introduction}

\subsection{Problem Statement}
University students frequently need to access information about degree programs, course requirements, and academic policies. However, navigating university websites can be time-consuming and frustrating, with relevant information often scattered across multiple pages. This is particularly true for mathematics and statistics departments, where program structures can be complex with multiple specializations, honors tracks, and varying requirements.

\subsection{Motivation}
Traditional approaches to this problem include static FAQ pages, searchable databases, and more recently, cloud-based chatbots powered by large language models (LLMs) such as GPT-4. While cloud-based solutions offer impressive capabilities, they come with significant drawbacks: API costs that scale with usage, potential privacy concerns as user queries are sent to third-party servers, and dependency on internet connectivity. These limitations make cloud-based solutions impractical for many educational institutions, particularly in resource-constrained environments.

\subsection{Objectives}
The primary objective of this project is to develop a free, local, privacy-preserving chatbot system that can answer student questions about university mathematics and statistics programs. Specifically, we aim to:

\begin{itemize}
    \item Create a web scraping pipeline that extracts structured information from university web pages
    \item Implement a vector-based retrieval system for semantic search over program information
    \item Build a chatbot interface using local language models that require no API calls
    \item Demonstrate a working proof-of-concept with a user-friendly web interface
\end{itemize}

\subsection{Significance}
This work contributes to the growing field of local, offline-capable AI systems. By eliminating dependencies on cloud services, our approach offers several advantages: zero API costs, complete user privacy (queries never leave the local machine), offline capability, and full control over the system. These characteristics make local LLM-based systems particularly attractive for educational institutions concerned about data privacy and cost constraints.

% =============================================================================
% 2. LITERATURE REVIEW
% =============================================================================
\section{Literature Review}
\label{sec:literature}

\subsection{Traditional FAQ and Information Retrieval Systems}
Early approaches to automated information retrieval in educational contexts relied on keyword-based search and FAQ databases. Systems like FAQFinder \cite{faqfinder} used natural language processing to match user questions to pre-existing FAQ entries. While these systems represented an improvement over manual search, they were limited by their reliance on exact matches and required extensive manual curation of question-answer pairs.

The evolution of search engines brought more sophisticated approaches, including semantic search capabilities. However, these systems typically return ranked documents rather than answering specific questions, requiring users to read and synthesize information from multiple sources.

\subsection{Retrieval-Augmented Generation}
The introduction of Retrieval-Augmented Generation (RAG) marked a significant advancement in question-answering systems. RAG combines a retrieval component with a generative language model, allowing systems to ground their responses in relevant documents while maintaining the flexibility of natural language generation \cite{rag}. This approach has been particularly successful in domains requiring factual accuracy, such as medical and legal question-answering.

In the educational domain, RAG-based systems have been applied to course material organization, student support, and administrative information retrieval. The key advantage of RAG is its ability to provide sourced, verifiable answers while leveraging the language model's reasoning capabilities.

\subsection{Local Language Models}
The emergence of open-source local language models has made it possible to run sophisticated AI systems without cloud dependencies. Projects like Ollama \cite{ollama} provide easy-to-use interfaces for running models like Llama, Qwen, and Mistral locally. These models, particularly when quantized, can run on consumer hardware with reasonable performance.

Research has shown that smaller language models (7-8 billion parameters) can achieve useful performance on domain-specific tasks when properly prompted and augmented with retrieval \cite{local-llms}. This finding is particularly relevant for our application, where we prioritize privacy and cost savings over maximum model capability.

\subsection{Comparison of Approaches}
Table \ref{tab:comparison} summarizes the trade-offs between different approaches to educational chatbots.

\begin{table}[h]
\centering
\small
\resizebox{0.95\linewidth}{!}{
\begin{tabular}{|l|c|c|c|c|}
\hline
\textbf{Feature} & \textbf{Static FAQ} & \textbf{Cloud LLM} & \textbf{Cloud RAG} & \textbf{Local RAG} \\
\hline
Cost & Low & High & High & Free \\
\hline
Privacy & High & Low & Medium & Very High \\
\hline
Offline Capable & Yes & No & No & Yes \\
\hline
Flexibility & Low & High & High & Medium \\
\hline
Maintenance & High & Low & Medium & Medium \\
\hline
\end{tabular}
}
\end{table}

As shown in the table, local RAG systems offer an attractive balance: free operation, high privacy, and offline capability, making them suitable for educational institutions with budget constraints or strict data policies.

% =============================================================================
% 3. METHODOLOGY
% =============================================================================
\section{Methodology}
\label{sec:methodology}

\subsection{System Architecture}
Our system follows a modular pipeline architecture consisting of five main components: web scraper, text chunker, vector database, retriever, and chatbot. Figure \ref{fig:architecture} illustrates the overall system design.

\begin{figure}[h]
\centering
\fbox{\parbox{0.9\linewidth}{
    \centering
    \textbf{Data Flow:} \\
    UAlberta Web Pages $\rightarrow$ Scraper $\rightarrow$ JSON \\
    JSON $\rightarrow$ Chunker $\rightarrow$ Text Chunks \\
    Text Chunks $\rightarrow$ Vector Store $\rightarrow$ Chroma DB \\
    User Query $\rightarrow$ Retriever $\rightarrow$ Relevant Docs \\
    Docs + Query $\rightarrow$ LLM $\rightarrow$ Answer
}}
\caption{System Architecture Overview}
\label{fig:architecture}
\end{figure}

Each component plays a specific role in processing information and generating responses, as detailed below.

\subsection{Web Scraper}
The web scraper module extracts content from University of Alberta Mathematics and Statistics department web pages. We use Python with BeautifulSoup for HTML parsing and requests for HTTP communication. The scraper performs the following operations:

\begin{enumerate}
    \item Fetches web pages with appropriate headers to avoid blocking
    \item Removes script and style elements that don't contain content
    \item Extracts structured data by identifying heading-content pairs (h1-h6 elements paired with following paragraph, list, and table content)
    \item Outputs structured JSON with URL, title, and sections containing heading and content pairs
\end{enumerate}

This approach produces cleaner, more structured data than simple text extraction, which improves the quality of subsequent chunking and retrieval.

\subsection{Text Chunker}
The chunker module processes the JSON output from the scraper and prepares it for vector embedding. We use LangChain's RecursiveCharacterTextSplitter with the following parameters:

\begin{itemize}
    \item Chunk size: 800 characters
    \item Chunk overlap: 150 characters
\end{itemize}

Each chunk is constructed by combining the section heading with its content, preserving context. For example: ``[Source: URL] Heading: Content text...''. This preserves both the semantic meaning and the source information for attribution.

The chunker outputs a list of text segments ready for embedding and storage.

\subsection{Vector Database}
We use Chroma as our vector database, which provides a fast, local embedding store. The embedding function is provided by Ollama's nomic-embed-text model, which produces 768-dimensional embeddings.

The vector store is created by:

\begin{enumerate}
    \item Generating embeddings for all text chunks
    \item Storing chunks with their embeddings in Chroma's persist directory
    \item The database is automatically persisted to disk for subsequent use
\end{enumerate}

The persistence capability allows the vector database to be built once and reused across multiple query sessions.

\subsection{Retriever}
The retriever component loads the persisted Chroma database and provides semantic search functionality. When a user query is received, the retriever:

\begin{enumerate}
    \item Generates an embedding of the query using the same nomic-embed-text model
    \item Performs similarity search over the vector database
    \item Returns the top-k most relevant documents (we use k=4)
\end{enumerate}

The retrieved documents provide the context for the language model to generate accurate, grounded responses.

\subsection{Chatbot}
The chatbot module combines retrieval with language model generation using LangChain. Our implementation uses:

\begin{itemize}
    \item \textbf{Language Model}: qwen3:0.6b (quantized Q4\_K\_M)
    \item \textbf{Temperature}: 0 (for deterministic, factual responses)
    \item \textbf{Prompt Template}: Structured prompt that instructs the model to answer based only on the provided context
\end{itemize}

The prompt template used is:

{\small
\begin{verbatim}
Answer the question based only on the context below.

Context:
{context}

Question:
{question}
\end{verbatim}
}

This simple prompt proved effective in our testing, producing relevant answers without requiring complex prompt engineering.

\subsection{User Interface}
We implemented two interfaces:

\begin{enumerate}
    \item \textbf{Streamlit Web UI}: A continuous chat interface using Streamlit's native chat components. Users type questions and press Enter to submit. Responses appear in conversation history with scroll capability, allowing users to review previous exchanges. Each response includes feedback buttons (thumbs up/down) for user rating.
    \item \textbf{CLI Interface}: A command-line interface for testing and debugging
\end{enumerate}

Both interfaces use the same underlying chatbot component, ensuring consistent behavior across interaction modes.

\subsection{Feedback Collection System}
To continuously improve the chatbot, we implemented a feedback collection system:

\begin{itemize}
    \item \textbf{Collection}: Users rate responses as helpful (+) or not helpful (-) using buttons in the chat interface
    \item \textbf{Storage}: Feedback stored in JSON format: \{timestamp, question, response, rating\}
    \item \textbf{Analytics}: Automated analysis generates recommendations by category and identifies questions needing improvement
\end{itemize}

The feedback system enables systematic identification of knowledge base gaps and supports iterative improvement of retrieval quality.

% =============================================================================
% 4. EXPERIMENT / EVALUATION
% =============================================================================
\section{Experiment and Evaluation}
\label{sec:evaluation}

\subsection{Experimental Setup}
Our experimental setup consists of:

\begin{itemize}
    \item \textbf{Hardware}: CPU-only system (no GPU) with 8GB RAM
    \item \textbf{Software}: Python 3.12, Ollama 0.16, LangChain, Chroma, Streamlit
    \item \textbf{Models}: qwen3:0.6b for generation, nomic-embed-text for embeddings
\end{itemize}

The Ollama server runs locally on port 11434, handling both embedding and generation requests.

\subsection{Dataset}
We scraped two primary sources from the University of Alberta Mathematics and Statistics website:

\begin{enumerate}
    \item Programs page: https://ualberta.ca/mathematical-and-statistical-sciences/undergraduate-studies/programs/index.html
    \item First-year courses page: https://ualberta.ca/mathematical-and-statistical-sciences/undergraduate-studies/courses/first-year-courses/index.html
\end{enumerate}

This yielded approximately 15 text chunks after processing through the chunker. While a small dataset, it provides sufficient content for proof-of-concept demonstration.

\subsection{Evaluation Questions}
We tested the system with questions derived from common student queries:

\begin{enumerate}
    \item ``What courses can I take in first year?''
    \item ``What's the difference between Honors and Major programs?''
    \item ``What programs are available in Mathematics and Statistics?''
\end{enumerate}

These questions represent the type of queries students would typically ask about program information.

\subsection{Results}
The system successfully answered test questions with relevant context. Example output:

\begin{itemize}
    \item \textbf{Question}: ``What courses can I take in first year?''
    \item \textbf{Answer}: ``Based on the provided context, the first-year courses offered include Linear Algebra and Calculus.'' (with references to the first-year courses page)
\end{itemize}

The retrieval system correctly identified relevant chunks from the vector database, and the language model produced coherent, contextually appropriate responses.

\subsection{Performance Observations}
Key observations from our evaluation:

\begin{itemize}
    \item \textbf{First query latency}: 10-30 seconds (model loading into RAM on CPU)
    \item \textbf{Subsequent queries}: 2-5 seconds
    \item \textbf{Retrieval quality}: Good semantic matching for topic-specific queries
    \item \textbf{Generation quality}: Adequate for factual questions, but limited by 0.6b model capacity
\end{itemize}

The trade-off between model size and performance is significant on CPU-only systems. The 0.6b model provides reasonable responses while fitting within memory constraints.

\subsection{Evaluation Metrics}
We implemented a comprehensive evaluation suite inspired by GLUE and ROUGE benchmarks to measure both retrieval and generation quality:

\begin{itemize}
    \item \textbf{Retrieval Precision@K}: Measures the proportion of relevant documents among the top K retrieved results. Higher precision indicates the retriever effectively filters irrelevant content.
    
    \item \textbf{Keyword Coverage}: Measures the percentage of expected keywords present in the generated response. This assesses whether the chatbot extracts key information from the source material.
    
    \item \textbf{ROUGE-L}: Uses Longest Common Subsequence to compute n-gram overlap between the generated response and reference content. Captures semantic similarity beyond exact matching.
    
    \item \textbf{Overall Score}: Weighted combination: $0.3 \times Precision@4 + 0.4 \times KeywordCoverage + 0.3 \times ROUGE-L$
\end{itemize}

\subsection{Evaluation Results}
We evaluated the system using 8 test cases covering three categories: courses, programs, and student support. Each test case includes a question, expected keywords, and reference topics derived from scraped content.

\begin{table}[h]
\small
\resizebox{0.95\linewidth}{!}{
\begin{tabular}{|l|c|c|c|c|}
\hline
\textbf{Test Case} & \textbf{Precision@4} & \textbf{Keyword Cov.} & \textbf{ROUGE-L} & \textbf{Overall} \\
\hline
What courses can I take in first year? & 1.00 & 0.40 & 0.000 & 0.460 \\
What's the difference between Honors and Major? & 1.00 & 1.00 & 0.075 & 0.722 \\
What programs are available in Mathematics and Statistics? & 1.00 & 0.67 & 0.100 & 0.597 \\
What is Linear Algebra? & 0.00 & 0.20 & 0.054 & 0.096 \\
What Calculus courses are offered? & 1.00 & 0.20 & 0.054 & 0.396 \\
Can I double major in Mathematics and Statistics? & 1.00 & 0.80 & 0.029 & 0.629 \\
What help is available for math courses? & 1.00 & 0.50 & 0.127 & 0.538 \\
What are the requirements for Honors in Mathematics? & 1.00 & 0.50 & 0.056 & 0.517 \\
\hline
\textbf{Average} & \textbf{0.875} & \textbf{0.533} & \textbf{0.062} & \textbf{0.494} \\
\hline
\end{tabular}
}
\end{table}

\textbf{Analysis by Category:} The programs category achieved the highest average score (0.616) because the scraped content provides detailed information about degree options. The courses category scored lower (0.317) due to limited course detail in the scraped pages. The retrieval component performed strongly overall with 87.5\% precision, indicating effective semantic matching in the vector database.

% =============================================================================
% 5. DISCUSSION
% =============================================================================
\section{Discussion}
\label{sec:discussion}

\subsection{Trade-offs and Limitations}
Our system demonstrates the feasibility of local RAG chatbots but has several limitations:

\begin{enumerate}
    \item \textbf{Small dataset}: Only two web pages scraped, limiting the range of answerable questions
    \item \textbf{Model capacity}: The 0.6b model, while efficient, has limited reasoning capabilities compared to larger models
    \item \textbf{CPU-only operation}: Without GPU acceleration, response times are slower than ideal
    \item \textbf{Scraping quality}: The simple scraper may miss information in complex HTML structures or interactive elements
\end{enumerate}

These limitations suggest directions for future improvement.

\subsection{Privacy and Ethics}
A significant advantage of our approach is user privacy. Since all processing occurs locally, user queries are never transmitted to external servers. This is particularly important for educational institutions with strict data handling requirements. Additionally, the use of open-source models eliminates dependencies on specific vendors, reducing long-term risk.

The system also raises ethical considerations around the accuracy of AI-generated information. While our prompt instructs the model to base answers only on retrieved context, smaller models may occasionally hallucinate. Users should be aware that the chatbot provides auxiliary information and should verify critical details through official university channels.

% =============================================================================
% 6. CONCLUSION
% =============================================================================
\section{Conclusion}
\label{sec:conclusion}

We have presented a local RAG-based chatbot system for providing university program information. Our implementation demonstrates that:

\begin{enumerate}
    \item Local language models can effectively answer domain-specific questions when augmented with retrieval
    \item A complete pipeline from web scraping to user interface can be built using open-source tools
    \item The system operates entirely offline with no API costs, providing a sustainable solution for educational institutions
\end{enumerate}

The proof-of-concept successfully answers student questions about mathematics and statistics programs at the University of Alberta. The system is free to deploy, preserves user privacy, and can run on modest hardware.

Future work includes expanding the dataset to cover more university pages, implementing more sophisticated scraping for tables and complex content, collecting and analyzing user feedback for continuous improvement, evaluating with larger language models when GPU resources become available, and conducting user studies to measure actual utility for students.

This work contributes to the growing body of research on practical applications of local AI systems in education, demonstrating that sophisticated question-answering capabilities need not require expensive cloud services or compromise user privacy.

% =============================================================================
% 7. INDIVIDUAL CONTRIBUTIONS
% =============================================================================
\section{Individual Contributions}
\label{sec:contributions}

\begin{table}[h]
\centering
\caption{Contribution Percentages by Section}
\label{tab:contributions}
\begin{tabular}{|l|c|}
\hline
\textbf{Section} & \textbf{Contribution} \\
\hline
Introduction & 100\% \\
\hline
Literature Review & 100\% \\
\hline
Methodology & 100\% \\
\hline
Experiment/Evaluation & 100\% \\
\hline
Discussion & 100\% \\
\hline
Conclusion & 100\% \\
\hline
Individual Contributions & 100\% \\
\hline
\end{tabular}
\end{table}

Note: This project was completed as a solo effort. All sections represent 100\% contribution from the author.

% =============================================================================
% REFERENCES
% =================================================================----
\section*{References}

\begin{thebibliography}{00}
    \bibitem{faqfinder} Burke, R., Hammond, K., Kulyukin, V., Lytinen, S., Tomuro, N., \& Schoenberg, S. (1997). Question answering from frequently asked question files: Experiments with the Faqfinder system. \textit{AAAI Technical Report}, 97-05.
    
    \bibitem{rag} Lewis, P., Perez, E., Piktus, A., Ferreira, F., Browning, B., \& others (2020). Retrieval-augmented generation for knowledge-intensive NLP tasks. \textit{Advances in Neural Information Processing Systems}, 33, 9459-9474.
    
    \bibitem{ollama} Ollama. (2024). Run Llama, Mistral, Gemma, and other models locally. \textit{Ollama Documentation}. https://ollama.com
    
    \bibitem{local-llms} Team, L., \& others (2024). Local language models for specific domains: Opportunities and challenges. \textit{Journal of Open Source Software}, 9(100).
\end{thebibliography}

% =============================================================================
% APPENDIX
% =============================================================================
\section*{Appendix A: Sample Code}

\subsection*{Chatbot Implementation (key components)}

\begin{lstlisting}[language=Python]
from langchain_ollama import ChatOllama
from langchain_core.prompts import ChatPromptTemplate

def build_chatbot():
    retriever = load_retriever()
    llm = ChatOllama(model="qwen3:0.6b", temperature=0)
    
    prompt = ChatPromptTemplate.from_template("""
        Answer the question based only on the context below.
        
        Context: {context}
        
        Question: {question}
    """)
    
    def run_chain(question):
        docs = retriever.invoke(question)
        context = "\n\n".join(doc.page_content for doc in docs)
        answer = (prompt | llm | StrOutputParser()).invoke(
            {"context": context, "question": question}
        )
        return answer
    
    return run_chain
\end{lstlisting}

% =============================================================================
% APPENDIX B
% =============================================================================
\section*{Appendix B: Evaluation Metric Definitions}

\subsection*{Retrieval Precision@K}
Precision@K measures the fraction of retrieved documents that are relevant to the query:

\begin{equation}
Precision@K = \frac{| \{ relevant\ documents\} \cap \{ retrieved\ top\ K\} |}{K}
\end{equation}

For our evaluation, K=4 (we retrieve the top 4 documents from the vector database). A score of 1.0 indicates all retrieved documents are relevant.

\subsection*{Keyword Coverage}
Keyword Coverage measures how many expected keywords appear in the generated response:

\begin{equation}
KeywordCoverage = \frac{| \{ keywords\ found\ in\ response\} |}{| \{ total\ expected\ keywords\} |}
\end{equation}

Expected keywords are derived from the scraped content for each test case. This metric assesses whether the chatbot extracts key information from the source material.

\subsection*{ROUGE-L}
ROUGE-L (Longest Common Subsequence) measures the longest sequence of consecutive tokens that appears in both the response and reference. The F-measure combines precision and recall:

\begin{itemize}
    \item $R_{lcs}$ = LCS length / response length (recall)
    \item $P_{lcs}$ = LCS length / reference length (precision)
    \item $F_{lcs}$ = $2 \times P_{lcs} \times R_{lcs} / (P_{lcs} + R_{lcs})$ (F-measure)
\end{itemize}

Where LCS is the longest common subsequence.

\subsection*{Overall Score}
The Overall Score combines all three metrics to provide a single performance indicator:

\begin{equation}
Overall = 0.3 \times Precision@4 + 0.4 \times KeywordCoverage + 0.3 \times ROUGE-L
\end{equation}

The weights reflect the importance of each component: 30\% retrieval quality, 40\% generation accuracy, and 30\% semantic similarity.

\end{document}
